\documentclass[10pt,a4paper]{amsart}
\usepackage[margin=19mm]{geometry}
\usepackage{amsmath,amssymb,mathtools}
\usepackage[scr=rsfs,cal=euler]{mathalfa}
\usepackage{xfrac}
\usepackage[unicode,hidelinks,bookmarks=false]{hyperref}
\newcommand{\ie}{i.e.\ }
\newcommand{\cf}{cf.\ }
\newcommand{\resp}{respectively\ }
\newcommand{\Subsection}{\subsection}
\providecommand{\nobreakdash}{\nobreak\hbox{-}}
\setlength{\emergencystretch}{2em}
\multlinegap0pt
\usepackage{amssymb}
\usepackage{algorithm}\usepackage{algpseudocode}
\usepackage{xcolor}
\usepackage[normalem]{ulem}
\newtheorem{lemma}{Lemma}
\newcommand{\X}{\mathcal{X}}
\DeclareMathOperator*{\divx}{div_{\X}}
\title{Static and Dynamic Approaches to Computing Barycenters of Probability Measures on Graphs}
\author{David Gentile, James M. Murphy}
\date{Source: arXiv:2603.26940v1 (CC BY 4.0)}
\begin{document}
\maketitle

\noindent\emph{Source and licence.} Static and Dynamic Approaches to Computing Barycenters of Probability Measures on Graphs. Version arXiv:2603.26940v1 (CC BY 4.0), distributed by arXiv under the Creative Commons Attribution 4.0 International licence, http://creativecommons.org/licenses/by/4.0/, as recorded in the arXiv licence metadata for this exact version and shown on the abstract page https://arxiv.org/abs/2603.26940v1. Full licence notice: \texttt{licenses/arxiv-2603.26940-CC-BY-4.0.txt}.\par\smallskip

\noindent\textit{Editorial scope.} Counted unit: the article's lemma lemma1 (source0.tex lines 592--598). The article's complete original proof of it is delivered with it (source0.tex lines 599--606, one \texttt{proof} environment of 8 lines). No mathematics is written, repaired or re-derived: the statement and the proof are the article's own lines, in the article's own notation, and no retained line is substituted or re-flowed.  No further source-local context is needed: every cross-reference of every retained line resolves inside the two delivered spans.  Nothing was omitted from any retained span, so the package carries no omission record.  No retained line contains a citation, so the wrapper carries no bibliography and no bibliography entry.  No retained line contains a figure environment, an image-inclusion command or drawing code, so no image is delivered and the package declares no asset dependency.  Editorial formatting only, in addition: an AMS \texttt{amsart} preamble that restates the article's own declarations and macros used by the retained lines, namely the environments \texttt{lemma} on separate counters, together with the standard packages those lines need.  The retained environments print in this document as Lemma 1 in order of appearance, whereas the article prints the same items with the numbering of its own full document.  The complete native source of the article as distributed in the arXiv e-print for this exact version is bundled unchanged as \texttt{sources/197314/source0.tex}.
\begin{lemma}
\label{lemma1}
Let \(\nabla\varphi \in T_{\nu}\mathcal{P}(\X)\). Then 
\begin{equation*}
\sum\limits_{x \in V}\divx \nabla\varphi(x)\pi(x) = 0.
\end{equation*}
\end{lemma}

\begin{proof}
    The claim follows from the reversibility of the pair \((Q, \pi)\). By definition of the graph gradient, \(\nabla\varphi\) is a skew-symmetric matrix. Therefore, we have 
    \begin{align*}
        \sum\limits_{x\in V}\divx\nabla\varphi(x)\pi(x) &= \sum\limits_{x\in V}\frac{1}{2}\sum\limits_{y \in V}Q(x,y)(\nabla\varphi(y,x) - \nabla\varphi(x,y))\pi(x), \\
      &=\sum\limits_{x,y\in V}\pi(x)Q(x,y)\nabla\varphi(y,x). 
    \end{align*}
Therefore we may interpret the graph divergence of \(\nabla\varphi\) as an elementwise sum of a certain matrix \(\hat{Q}\odot \nabla\varphi\), where \(\hat{Q}\) is the matrix \(\hat{Q} = [\pi(x)Q(x,y)]_{x,y \in \X}\) and \(\odot\) is the Hadamard product for matrices (i.e., elementwise multiplication). By reversibility of the pair \((Q, \pi)\), the matrix \(\hat{Q}\) is symmetric, so its Hadamard product with \(\nabla\varphi\) is skew-symmetric. Observing that the elementwise sum of any skew-symmetric matrix is \(0\) finishes the proof. 
\end{proof}

\end{document}
